\clearpage
\subsection{Entregables, estándares utilizados y justificación [Pendientes: 1. acordar criterios cuantitativos de aceptación]}\label{sec:entregables}
Los entregables se identificarán por versión y se relacionarán con los objetivos y las evidencias de cumplimiento, de modo que puedan ser revisados durante las fases de la sección~\ref{sec:proceso}. La tabla~\ref{tab:entregables} concentra los productos comprometidos y el propósito de los estándares seleccionados. Su utilización será una adaptación académica explícita al alcance del proyecto, por lo que citar una norma no implica certificación ni demuestra conformidad completa. Asimismo, los estándares de ingeniería y calidad apoyarán la organización de la evidencia, mientras que la validez de datos, simulaciones y agregación deberá sustentarse mediante comprobaciones técnicas específicas.

{\small\setlength{\tabcolsep}{5pt}\renewcommand{\arraystretch}{1.18}
\begin{longtable}{|>{\raggedright\arraybackslash}p{0.26\linewidth}|>{\raggedright\arraybackslash}p{0.28\linewidth}|>{\raggedright\arraybackslash}p{0.36\linewidth}|}
\caption{Entregables y aplicación de estándares.}\label{tab:entregables}\\\hline
\textbf{Entregable} & \textbf{Estándares asociados} & \textbf{Justificación}\\\hline\endfirsthead
\hline\textbf{Entregable} & \textbf{Estándares asociados} & \textbf{Justificación}\\\hline\endhead
E1. Especificación de requisitos y matriz de trazabilidad. & ISO/IEC/IEEE 29148:2018~\cite{iso29148}. & Relacionar necesidades, alcance, requisitos verificables y evidencia, conservando la procedencia de las necesidades y las exclusiones acordadas.\\\hline
E2. Documento de arquitectura y diseño. & OMG UML 2.5.1~\cite{uml251} e ISO/IEC/IEEE 12207:2026~\cite{iso12207}. & Describir componentes, interfaces y flujo de información, vinculando el diseño con el desarrollo e integración de la solución.\\\hline
E3. Paquete reproducible de configuración WFDS y preparación de datos. & ISO/IEC/IEEE 12207:2026 e IEEE 730:2026~\cite{ieee730}. & Identificar versiones, entradas, unidades, procedencia y evidencia de revisión de los productos técnicos.\\\hline
E4. Biblioteca de escenarios y catálogo de metadatos. & ISO/IEC/IEEE 12207:2026 e IEEE 730:2026. & Conservar la relación entre cada resultado, sus condiciones y su configuración, para permitir consulta e identificar casos incompletos o incompatibles.\\\hline
E5. Código fuente y versión ejecutable de la aplicación de escritorio. & ISO/IEC/IEEE 12207:2026 y OMG UML 2.5.1. & Mantener correspondencia entre diseño, implementación e integración de los componentes de consulta y visualización.\\\hline
E6. Especificaciones, registros e informe de pruebas y evaluación. & ISO/IEC/IEEE 29119-2, 29119-3 y 29119-4:2021~\cite{iso291192,iso291193,iso291194}, e IEEE 730:2026. & Organizar procesos, documentación y técnicas de prueba, registrando condiciones, resultados y reevaluaciones frente a criterios explícitos.\\\hline
E7. Manual de instalación y uso e informe final del proyecto. & ISO/IEC/IEEE 12207:2026 e IEEE 730:2026. & Acompañar la entrega con instrucciones reproducibles, dependencias, límites conocidos y síntesis de la evidencia obtenida.\\\hline
\end{longtable}}
